\documentclass[10pt,a4paper]{amsart}
\usepackage[margin=19mm]{geometry}
\usepackage{amsmath,amssymb,mathtools}
\usepackage[scr=rsfs,cal=euler]{mathalfa}
\usepackage{xfrac}
\usepackage[unicode,hidelinks,bookmarks=false]{hyperref}
\newcommand{\ie}{i.e.\ }
\newcommand{\cf}{cf.\ }
\newcommand{\resp}{respectively\ }
\newcommand{\Subsection}{\subsection}
\providecommand{\nobreakdash}{\nobreak\hbox{-}}
\setlength{\emergencystretch}{2em}
\multlinegap0pt
\usepackage{amssymb}
\usepackage[shortlabels]{enumitem}
\usepackage{mathtools}
\newtheorem{theorem}{Theorem}
\newtheorem{lemma}{Lemma}
\def\N{{\mathbb  N}}
\def\Q{\mathbb{Q}}
\newcommand{\T}{\mathbb T}
\def\Z{{\mathbb  Z}}
\newcommand{\dist}{\operatorname{dist}}
\newcommand{\nnorm}[1]{\left\|#1\right\|}
\title{Polynomial extensions of Raimi's theorem}
\author{Norbert Hegyvari, Janos Pach, Thang Pham}
\date{Source: arXiv:2511.06650v3 (CC BY 4.0)}
\begin{document}
\maketitle

\noindent\emph{Source and licence.} Polynomial extensions of Raimi's theorem. Version arXiv:2511.06650v3 (CC BY 4.0), distributed by arXiv under the Creative Commons Attribution 4.0 International licence, http://creativecommons.org/licenses/by/4.0/, as recorded in the arXiv licence metadata for this exact version and shown on the abstract page https://arxiv.org/abs/2511.06650v3. Full licence notice: \texttt{licenses/arxiv-2511.06650-CC-BY-4.0.txt}.\par\smallskip

\noindent\textit{Editorial scope.} Counted unit: the article's theorem thmpolynomial-main (source0.tex lines 248--261). The article's complete original proof of it is delivered with it (source0.tex lines 548--599, one \texttt{proof} environment of 52 lines, which the article places away from the statement). No mathematics is written, repaired or re-derived: the statement and the proof are the article's own lines, in the article's own notation, and no retained line is substituted or re-flowed.  Retained uncounted as the source-local context the proof needs: lines 385--388; lines 450--455; lines 523--525.  Nothing was omitted from any retained span, so the package carries no omission record.  No retained line contains a citation, so the wrapper carries no bibliography and no bibliography entry.  No retained line contains a figure environment, an image-inclusion command or drawing code, so no image is delivered and the package declares no asset dependency.  Editorial formatting only, in addition: an AMS \texttt{amsart} preamble that restates the article's own declarations and macros used by the retained lines, namely the environments \texttt{theorem}, \texttt{lemma} on separate counters, together with the standard packages those lines need.  The retained environments print in this document as Theorem 1, Lemma 1 in order of appearance, whereas the article prints the same items with the numbering of its own full document.  The complete native source of the article as distributed in the arXiv e-print for this exact version is bundled unchanged as \texttt{sources/188751/source0.tex}.
\begin{lemma}\label{lm2.1}
Let $\bigcup_{i=1}^tF_i$ be any partition (coloring) of $\N^k$. There exist an $F_m$ and an interval $(x,y)\subseteq [0,1)$ such that  $\varphi(F_m)$ is everywhere dense in $(x, y)$. 
  
\end{lemma}

\begin{lemma}\label{lm2.3}
Let $A\subset\T$ be an interval and $\sigma,\gamma\in\T$ with $\nnorm{\gamma-\sigma}<\varepsilon$. Then
\[
(A+\sigma)_{-\varepsilon}\ \subset\ A+\gamma\ \subset\ (A+\sigma)_{+\varepsilon}.
\]
\end{lemma}

\begin{theorem}\label{lm2.2'}
Let $\beta \notin \Q$. Let $P^{(1)},\dots,P^{(f)}\in\Z[x]$ be non-constant polynomials satisfying $P^{(1)}(0)=\cdots=P^{(f)}(0)=0$. For every $\varepsilon>0$, there exists a set $X_{\varepsilon}\subset \mathbb{N}$ of positive lower density such that for each fixed $h\in X_{\varepsilon}$,  for every $i\in \{1,\dots,f\}$, we have $\nnorm{\{\beta P^{(i)}(h)\}}<\varepsilon$. 
\end{theorem}

\begin{theorem}\label{thmpolynomial-main}
Let $r, t, k, f\in \mathbb{N}$.
Let $P^{(1)},\dots,P^{(f)}\in\Z[x]$ be non-constant polynomials with the properties that $P^{(1)}(0)=\cdots=P^{(f)}(0)=0$ and the leading coefficients are positive. 

There exists a partition
$\bigcup_{i=1}^rE_i$ of $\mathbb{N}^k$ such that for every finite coloring of $\mathbb{N}^k$ with $t$ colors,
$
\mathbb{N}^k = \bigcup_{j=1}^t F_j,
$ there exist a single color class $F_m$, $1\le m\le t$, an element $x_0\in \mathbb{N}$, and a set $H\subset\mathbb{N}$ of positive lower density with the property that for every $h\in H$ and every $j\in\{1,\dots,f\}$, the set
\[
\big(F_m+(x_0+P^{(j)}(h))\mathbf{1}_k\big)\cap E_i
\]
has infinitely many elements.
\end{theorem}

\begin{proof}[Proof of Theorem \ref{thmpolynomial-main}]
By Lemma~\ref{lm2.1}, choose $m$ and an interval $J=(x,y)\subset(0,1)$ such that $\varphi(F_m)$ is dense in $J$.
Write $\delta:=y-x>0$.

Choose a band $S_{j^\ast}$ with width $|S_{j^\ast}|=2^{-(j^\ast+1)}\le \delta/3$.
Pick $x_0\in\mathbb N$ and put $\sigma:=\{x_0\beta\}$, so that
\[
  J+\sigma\ \supset\ S_{j^\ast}\quad\text{mod }1.
\]
Choose $\varepsilon>0$ so small that
\[
  S_{j^\ast}\ \subset\ (J+\sigma)_{-\varepsilon}
  :=\{u\in J+\sigma:\ \dist(u,\partial(J+\sigma))>\varepsilon\}.
\]

It follows from Theorem \ref{lm2.2'} that there exists a set $X_{\varepsilon}\subset \mathbb{N}$ of positive lower density such that for any $h\in X_{\varepsilon}$ one has
\[
  \|\{\beta P^{(j)}(h)\}\|<\varepsilon \quad\text{for all } j=1,\dots,f.
\]
%{\color{blue}The lower density of $X_{\varepsilon}$ depends on $\varepsilon$ and on the polynomials, that is, on  $J, \sigma$, and on the polynomials.}

Moreover, since the leading coefficient of each $P^{(j)}$ is positive,
we may choose $h$ large enough to make sure that $P^{(j)}(h)\ge 0$ for all $j$. The set of $h$ satisfying this condition still has positive lower density.

For each fixed $j$, set
\[
  \gamma_j':=\sigma+ \{\beta P^{(j)}(h)\}\mod 1,
\]
Then $\|\gamma_j'-\sigma\|=\|\{\beta P^{(j)}(h)\}\|<\varepsilon$, and by Lemma~\ref{lm2.3},
\[
  (J+\sigma)_{-\varepsilon}\ \subset\ J+\gamma_j'\ \subset\ (J+\sigma)_{+\varepsilon}.
\]
In particular,
\[
  J+\gamma_j'\ \supset\ (J+\sigma)_{-\varepsilon}\ \supset\ S_{j^\ast}.
\]

Since $P^{(j)}(h)\ge 0$, the translate $F_m+(x_0+P^{(j)}(h))\mathbf{1}_k$ lies in $\mathbb N^k$, and
\[
  \varphi\!\big(\mathbf{x}+(x_0+P^{(j)}(h))\mathbf{1}_k\big)
  \;=\;\varphi(\mathbf{x})+\gamma_j' \mod 1.
\]
Because $\varphi(F_m)$ is dense in $J$, it follows that $\varphi(F_m+(x_0+P^{(j)}(h))\mathbf{1}_k)$
is dense in $J+\gamma_j'$, hence intersects each tile $T_{j^\ast,i}\subset S_{j^\ast}$ infinitely many times.
By the definition of the color classes $E_i$, we conclude that
\[
  \big(F_m+(x_0+P^{(j)}(h))\mathbf{1}_k\big)\cap E_i
  \quad\text{is infinite for every } i=1,\dots,r.
\]
This holds for every $j=1,\dots,f$ at the same $h$.
Since there are infinitely many such $h$,
the theorem follows. \end{proof}

\end{document}
